% Round 18. Application overlay; authored after the RI18 independent review.
\section{All finite fields for the complete product and ambient domains}
\label{r18ru:section}

This short overlay records the extension range supplied by RI18 after
its independent review.  It leaves the earlier HP18, NF18, GF18,
NG18, GH18 and CA18 sources and their originally stated scopes intact.
No new-cycle assertion is obtained by presuming that cycles descend.

\begin{corollary}[No new obstruction directions under totally ramified inertia]
\label{r18ru:quotient}
Let $X/K$ have the supplied native reduced strict semistable model
of RI18, and let $L/K$ be finite, totally ramified and Galois of
degree $d$.  For any $1\le r\le n$, put
\[
 \mathcal Q_K^r(X)=
 \operatorname{CH}^r(X)_{\mathbf Q}/\ker\mathfrak o_{X,K}^r.
\]
The ordinary Chow pullback induces an isomorphism
\[
 p^*:\mathcal Q_K^r(X)\xrightarrow{\ \sim\ }
            \mathcal Q_L^r(X_L),
       \qquad (p^*)^{-1}=d^{-1}p_* .
 \tag{RU18.0}\label{r18ru:quotient-iso}
\]
The same statement holds with each Chow group restricted to its
geometrically homologically trivial subgroup and with the
corresponding intersection of the obstruction kernel.
This is an assertion about the obstruction quotients, not about
Galois triviality of the entire Chow group.
\end{corollary}

\begin{proof}
Frozen FG16 gives both preservation of the existence kernel under
$p^*$ and injectivity on its quotient; VT17 identifies that existence
kernel with the indicated obstruction kernel.  RI18's equality
$\nu_{g z}=\nu_z$ for every rational Chow class shows that
$g z-z$ belongs to the upstairs kernel for every $g$.
Consequently the Galois identity $p^*p_*z=\sum_g gz$ gives
\[
                [p^*p_*z]=d[z]\quad\hbox{in }\mathcal Q_L^r(X_L).
\]
If $z$ lies in the upstairs kernel, the same identity and the exact
period detector imply that $p^*p_*z$ lies in that kernel.  FG16's
finite-field descent then implies that $p_*z$ lies in the downstairs
kernel.  Thus $p_*$ really does induce a map on these quotients;
this is not being presumed from an unspecified proper actual
pushforward.  The Chow degree identity $p_*p^*=d$ and the displayed
upstairs identity show that $d^{-1}p_*$ and $p^*$ are mutual inverses.
Both Chow operations preserve geometric homological triviality,
so the identical argument restricts to that subgroup.
\end{proof}

\begin{corollary}[Complete ambient domains over every finite valued field]
\label{r18ru:local}
Take all the native hypotheses of CA18: a smooth good-reduction
complete-intersection ambient variety of relative dimension $n+1$,
$n\ge4$, its specified geometric SNC factor divisors, every actual
complete-intersection stratum with the stated connectedness and
zero-dimensional conventions, the transverse degree-$D$ form $Q$,
and the exponent-one smoothing with smooth geometrically integral
generic fiber $X/K$.  Then for every finite valued extension $L/K$
and every specified geometrically homologically trivial rational
cycle $Z$ on $X_L$, in every codimension, there is a strong Green
form on the complement of exactly $|Z|$ against every compact
ordinary actual test, with the absolute value extending the fixed
one on $K$.
This includes every HP18/NF18/GH18 model and the type-$(2,5)$
fourfolds of CA18.
\end{corollary}

\begin{proof}
CA18 constructs an actual regular projective reduced strict
semistable native model.  Its local full-G theorem is explicitly
reapplied to all new cycles after every finite unramified extension:
such base change preserves every geometric stratum condition and
the exponent-one local equation.  Hence, for each codimension,
this supplied model satisfies exactly the hypothesis of
RI18's all-finite-field corollary.
That corollary handles a specified cycle over an arbitrary $L$
by a Galois closure, a fresh unramified-field theorem, inertial
norm detection, and full finite actual trace.  It supplies the
claimed form for that original cycle and its precise support.
All field degrees are the complete degrees, and the absolute value
is never renormalized.  The specializations named in the statement
are cases of the same CA18 hypotheses; their cycle groups are not
replaced by images of earlier-field Chow groups.
\end{proof}

\begin{corollary}[A finite-place criterion stable under every number-field extension]
\label{r18ru:global-criterion}
Let $F$ be a number field and $X/F$ a smooth projective geometrically
integral variety of dimension $n\ge4$.  Suppose that for each finite
place $v$ of $F$ there is a finite valued extension $L_v/F_v$ over
which $X_{L_v}$ has one of the following supplied local descriptions:
\begin{enumerate}
\item a smooth projective good-reduction model;
\item an ordinary-quadratic model in the precise frozen ON17/GP17
      domain, including its allowed finite thickness and auxiliary
      native splitting;
\item one of the complete-intersection ambient product models of
      CA18, with all hypotheses of Corollary~\ref{r18ru:local}.
\end{enumerate}
Then, for every finite extension $E/F$, every finite place $w$ of
$E$, and every specified geometrically homologically trivial
rational cycle on $X_{E_w}$, strong G holds in every codimension.
Consequently the frozen PH16 construction and GA17 model-lift theorem
give, for every $1\le r\le n$, the complete complementary-degree
height pairing
\[
 \operatorname{CH}^r(X_E)_{\mathbf Q,\mathrm{geom\,hom}}
 \times
 \operatorname{CH}^{n+1-r}(X_E)_{\mathbf Q,\mathrm{geom\,hom}}
                         \longrightarrow\mathbf R.
 \tag{RU18.1}\label{r18ru:height}
\]
\end{corollary}

\begin{proof}
Fix $w\mid v$ and an arbitrary specified local cycle $Z$ over
$E_w$.  Embed the finite local fields in a common algebraic closure
of $F_v$ and take their compositum $M=L_vE_w$.
In the good-reduction case, smooth proper base change and the frozen
model detector apply anew to every homological cycle over $M$.
In the ordinary-quadratic case, the stated frozen theorem is reapplied
over $M$; its allowed thickness changes by the actual ramification
index, and its finite native auxiliary extensions are included in
its theorem.  In the ambient-product case, apply
Corollary~\ref{r18ru:local} to the finite extension $M/L_v$ and
the pullback of the present $Z$ to $M$.

Thus in each case that pulled specified cycle has strong G over
$M$.  The frozen finite actual trace descends it to $E_w$ with
exactly the degree $[M:E_w]$, or with the corresponding complete
composite degree if the frozen ordinary-quadratic theorem used a
further auxiliary field.  The original absolute value on $E_w$
and the precise support of $Z$ are retained.  Since $Z$ was arbitrary,
this proves the full local domain at every $w$, including classes
which are not defined over either $F_v$ or $L_v$.

PH16 now has its full local hypotheses in both complementary
codimensions on the entire displayed geometric homological groups.
GA17 supplies the required admissible model lifts on this established
domain.  Their frozen construction gives
\eqref{r18ru:height}.  It is not being extended from a smaller
already-defined height subgroup by choosing values on a quotient.
\end{proof}

\begin{corollary}[The unchanged GF18 and NG18 dense families]
\label{r18ru:existing-global-families}
Every member of the complete all-dimensional GF18/NG18 families
as upgraded in GH18 has the height pairings
\eqref{r18ru:height} after every number-field extension of
$\mathbf Q$.  The condition ``unramified at the distinguished prime
$q$'' in the earlier source statements is removed for these same
members.  Their coefficient conditions, density assertions and
monodromy ranks are unchanged.
\end{corollary}

\begin{proof}
GF18/NG18 establish the product description at $q$ and the frozen
good or ordinary-quadratic description at every other finite place.
Thus Corollary~\ref{r18ru:global-criterion} applies with $F=\mathbf Q$;
no new arithmetic spreading claim is necessary.
\end{proof}

\paragraph{Exact limits.}
A single local CA18 example is not, by itself, a global height
example: the finite-place hypothesis of
Corollary~\ref{r18ru:global-criterion} must be checked at every
other place.  A separate arithmetic construction of such examples
may invoke that criterion after its placewise verification.
These assertions cover the stated full cycle domains; they assert
neither every smooth complete intersection nor every hypersurface,
and they do not resolve general G or full BB heights for an arbitrary
smooth projective variety.  In particular RI18 has not removed its
need for a fresh full theorem over all unramified extensions.

\paragraph{Dependency and review status.}
The new input of this overlay is RI18's independently reviewed norm
detection theorem.  The all-unramified starting theorem is GH18 or
CA18, with its separately bound hypotheses and review; the full
height construction is PH16/GA17.  The global specializations use
the previously reviewed placewise GF18/NG18 construction.  This
application overlay has no literature-priority claim and is to be
bound separately by the final review/register.
